\AtBeginDocument{%
\begingroup\pagestyle{empty}\raggedright\parindent0pt
%\vspace*{-90pt}
%\begin{changemargin}{-60pt}{-60pt}
%{\fontsize{70}{1000}\Cobraarisca{
%\mbox{\noindent\rlap{\enspace C C C C C}{C C C C C C}}
%\mbox{\noindent\rlap{\enspace C C C C C}{C C C C C C}}
%\mbox{\noindent\rlap{\enspace C C C C C}{C C C C C C}}
%\mbox{\noindent\rlap{\enspace C C C C C}{C C C C C C}}
%\vspace{-18pt}{
%\mbox{\noindent\rlap{\enspace C C C C C}{C C C C C C}}}
%\mbox{\noindent\rlap{\enspace C C C l C}{C C C C C C}}
%\mbox{\noindent\rlap{\enspace C C C C C}{C C C C C C}}
%\mbox{\noindent\rlap{\enspace C C C C C}{C C C C C C}} \enlargethispage{\textheight}
%\mbox{\noindent\rlap{\enspace C C C C C}{C C C C C C}}
%\mbox{\noindent\rlap{\enspace C C C C C}{C C C C C C}}
%\mbox{\noindent\rlap{\enspace C C C C C}{C C C C C C}}
%}}
%\end{changemargin}
\begin{flushright}
\begin{vplace}[30]
``(\ldots{}) Um homem toma posse de si mesmo por meio de lampejos, e muitas vezes quando toma posse de si não se encontra nem se alcança. (\ldots{})''

\vspace*{0.2cm}

A. Artaud,
\textit{Carta para Jacques Rivière} em 25 de maio de 1924
\end{vplace}
\end{flushright}

%\LARGE{\autor}
%\vfill
\clearpage

%% Créditos ------------------------------------------------------
\raggedright

%{\Formular{\normalsize{N-1 edições \hspace{4.6cm} Coleção
%
%\vspace*{-5mm}
%
%\& Hedra \hspace{5cm} Lampejos}}}

%\bigskip
\linhalayout{Coleção Lampejos}{}
{\tiny{©n-1 edições 2020 / Hedra}}

\vspace*{0.5cm}

\linhalayout{\emph{Semente de crápula. Conselhos aos educadores}}{}
\linhalayout{\emph{que gostariam de cultivá-la}}{}
\linhalayout{Fernand Deligny}{}
\linha{título original}{\titulooriginal}
{\tiny{©n-1 edições 2020}}

\vspace*{1cm}

\linha{tradução©}{\copyrighttraducao}
\linhalayout{coordenação editorial}{Peter Pál Pelbart e Ricardo Muniz Fernandes}
\linhalayout{direção de arte}{Ricardo Muniz Fernandes}
\linhalayout{revisão}{Pedro Taam}
\linhalayout{revisão técnica da tradução}{Cícero Oliveira e Juliana Jardim}
\linhalayout{projeto da coleção/capa}{Lucas Kröeff}
\linhalayout{projeto}{Juliana Jardim e Ensaios ignorantes (ensaiosignorantes.com)}
%\linhalayout{ilustração/alfabeto}{Waldomiro Mugrelise}
\linhalayout{desenhos}{Fernand Deligny}
\linha{coedição}{\coedicao}
\linha{assistência editorial}{\assistencia}
\linha{ISBN}{\ISBN}\smallskip
\linha{revisão}{\revisao}
\linha{preparação}{\preparacao}
\linha{iconografia}{\iconografia}
\linhalayout{colaboração}{Maxime Godard e Sandra Alvarez de Toledo}
\linha{imagem da capa}{\imagemcapa}\smallskip
\begingroup\tiny
%\ifdef{\conselho}{\conselho}{\relax}
\par\endgroup\bigskip
%
\begingroup \tiny
%
\textit{Grafia atualizada segundo o Acordo Ortográfico da Língua\\
Portuguesa de 1990, em vigor no Brasil desde 2009.}\\

%
\begin{vplace}[30]
\vfill\textit{Direitos reservados em l\'ingua\\ portuguesa somente para o Brasil}\\\medskip
%%
\textsc{n-1 edições}\\
R.~Fradique Coutinho, 1139\\
05416--011 São Paulo \textsc{sp} Brasil\\
\end{vplace}

\endgroup

%{\scriptsize\fakereceipt{
%A CONTRACULTURA,\\
%ENTRE A CURTIÇÃO\\
%E O EXPERIMENTAL

%\vspace{1cm}

%CELSO FAVARETTO

%\vspace{2cm}

%DESENHOS\\
%\vspace{0.5cm}
%WALDOMIRO\\
%\hspace{1.7cm}MUGRELISE\\

%\vspace{1cm}

%COLEÇÃO LAMPEJOS

%\vspace{0.5cm}

%PRIMEIRA EDIÇÃO

%\vspace{0.5cm}

%COPYRIGHT: HEDRA \& N-1\\
%EDIÇÃO: RICARDO MUNIZ\\
%PROJETO DA COLEÇÃO: L.KRÖEFF\\
%TIRAGEM: 1\textsuperscript{A TIRAGEM} 300UN\\
%ISBN: XXXXXXXXXXXXXXXXXXXXX\\

%\vspace{1cm}

%PUBLICADO POR\\
%HEDRA \& N-1\\
%SÃO PAULO, 2019\\
%IMPRESSA SOB DEMANDA\\

%\vspace{0.5cm}

%RUA FRADIQUE COUTINHO, 1139\\
%05416-011 SÃO PAULO, SP BRASIL
%}}
\pagebreak

%\vspace*{-80pt}
\begin{changemargin}{-100pt}{-60pt}
{\Formular{

\mbox{\textbf{Fernand Deligny} \hspace{1.8cm}\textbf{Fernand Deligny} \hspace{1.8cm}\textbf{Fernand Deligny}}\vspace*{-3pt}
\small{\mbox{Semente de crápula \hspace{1.6cm}Semente de crápula \hspace{1.65cm}Semente de crápula}\vspace*{-3pt}}
\scriptsize{\mbox{Conselhos aos educadores \hspace{1.2cm}Conselhos aos educadores \hspace{1.3cm}Conselhos aos educadores}\vspace*{-7pt}
\mbox{que gostariam de cultivá-la \hspace{1.2cm}que gostariam de cultivá-la \hspace{1.3cm}que gostariam de cultivá-la}}
%\mbox{\textbf{Fernand Deligny} \hspace{0.95cm}\textbf{Fernand Deligny} \hspace{0.95cm}\textbf{Fernand Deligny} \hspace{0.95cm}\textbf{Fernand Deligny}}\vspace*{-3pt}
%\mbox{\small{Semente de crápula:} \hspace{.4cm}Semente de crápula: \hspace{.4cm}Semente de crápula: \hspace{.4cm}Semente de crápula:}\vspace*{-3pt}
}}
\end{changemargin}
%% Front ---------------------------------------------------------
% Titulo


%{\LARGE{\autor} \par}%\vspace{1.5ex}}
%\vspace{6cm} %9.3cm
%\ifdef{\organizador}{{\small {\organizador} (\textit{organização})} \par}{}
%\ifdef{\introdutor}{{\small {\introdutor} (\textit{prefácio})} \par}{}
%\ifdef{\tradutor}{{\small {\tradutor} (\textit{tradução})}\par}{}\vspace{8.5mm}

%{{\footnotesize{} \ifdef{\numeroedicao}{\numeroedicao}{2}ª edição} \par}
%logos
\vfill
%\normalsize
%\ifdef{\logo}{\IfFileExists{\logo}{\hfill\includegraphics[width=3cm]{\logo}\hfill\logoum{}\\ São Paulo\_\the\year}}{\logoum\break{} São Paulo\_\the\year}
%\includegraphics[width=.4\textwidth,trim=0 0 25 0]{logo.jpg}\\\smallskip
\par\clearpage\endgroup
% Resumo -------------------------------------------------------
\begingroup \footnotesize \parindent0pt \parskip 5pt \thispagestyle{empty} \vspace{1cm}%\textheight}\mbox{} \vfill
\movetoevenpage
\raggedright
{\Formular{\textit{O livro como imagem do mundo é de toda maneira uma ideia insípida. Na verdade não basta dizer Viva o múltiplo, grito de resto difícil de emitir. Nenhuma habilidade tipográfica, lexical ou mesmo sintática será suficiente para fazê-lo ouvir. É preciso fazer o múltiplo, não acrescentando sempre uma dimensão superior, mas, ao contrário, da maneira mais simples, com força de sobriedade, no nível das dimensões de que se dispõe, sempre n-1 (é somente assim que o uno faz parte do múltiplo, estando sempre subtraído dele). Subtrair o único da multiplicidade a ser constituída; escrever a n-1.}}}

\vspace*{4pt}

{\Formular{\textit{Gilles Deleuze e Félix Guattari}}}
\movetooddpage
\baselineskip=.92\baselineskip
\IfFileExists{PRETAS.tex}{\raggedright

{\Formular{\textbf{FERNAND DELIGNY} \textit{nasceu no norte da França em 1913. Conhecido
como educador, preferia autonomear"-se poeta e etólogo: “Meu projeto
era escrever”. A escrita é nele uma atividade existencial, o
laboratório permanente da sua prática. Dedica 50 anos da vida a crianças
inadaptadas, delinquentes, psicóticos e autistas. Etólogo designa o meio
que ele reinventa, em todas as suas circunstâncias, para tentar dar a
essas crianças a oportunidade de sobreviver em uma comunidade que exclui
ou normaliza.}

\textit{Deligny publicou 15 livros, mais de cem artigos, fez filmes, fotografou e
escreveu milhares de páginas inéditas: novelas inacabadas, roteiros,
contos, ensaios infinitamente recomeçados sobre o humano, a linguagem, o
psiquiátrico, a rede, o aracnídeo; esses dois últimos termos são marcas
de seu glossário. Dois de seus livros, {\slsc{Os vagabundos eficazes}} e
{\slsc{O aracniano}}, estão publicados pela n-1.}
}}




}{% 
\ifdef{\resumo}{\resumo\par}{}
\ifdef{\sobreobra}{\sobreobra}{}
\ifdef{\sobreautor}{\mbox{}\vspace{4pt}\newline\sobreautor}{}
\ifdef{\sobretradutor}{\newline\sobretradutor}{\relax}
\ifdef{\sobreorganizador}{\vspace{4pt}\newline\sobreorganizador}{\relax}\par}
\thispagestyle{empty} \endgroup
\ifdefvoid{\sobreautor}{}{\pagebreak\ifodd\thepage\paginabranca\fi}
% Sumário -------------------------------------------------------

%\sumario{}
%\IfFileExists{INTRO.tex}{\include{INTRO}}

%\IfFileExists{TEXTO.tex}{\mbox{}\include{TEXTO}}
%\part[{{\def\break{}\titulo}}]{\titulo}
} % fim do AtBeginDocument

% Finais -------------------------------------------------------
\AtEndDocument{%
  %\publicidade
%
\pagebreak\ifodd\thepage\paginabranca\fi
\thispagestyle{empty}
\movetooddpage

\vspace*{\fill}
{\footnotesize\Formular\raggedright
A ideia desta coleção Lampejos foi criar, para cada capa, um alfabeto diferente desenhado pelo artista Waldomiro Mugrelise. Entremear a singularidade dos textos de cada autor à invenção gráfica de um outro léxico e outra sintaxe.

\bigskip

“Todos os viajantes confirmaram: transformar o teclado do computador em mecanismo de fazer desenhos é a melhor solução para este projeto. A invenção de um dispositivo composicional além do léxico, quero dizer, anterior ao léxico, fará o leitor percorrer léguas de insensatas cacofonias, de confusões verbais e repetições que correspondem a idioma algum, por dialetal ou rudimentar que seja. A incoerência (inocorrência?) da palavra resulta em potencialidade gráfica infinita, um campo ilimitado para o desenho. Lucas compõe as capas a partir da tipologia fornecida por Waldomiro. Eu me visto de Waldomiro, diz ele. Ser meio para nenhum fim. As linhas caóticas da mão são capturadas e organizadas em um sistema que produz composições que o artista nunca criaria. Imagem é texto, como bem sabemos. Os livros, por diversos que sejam, constam de elementos iguais: o espaço, o ponto, a vírgula, as letras do alfabeto.” 

\hfill{}{\italic{Leopardo Feline}}
}
\vspace*{\fill}

\pagebreak\ifodd\thepage\paginabranca\fi
\thispagestyle{empty}
{\raggedright
\includegraphics[width=.4\textwidth]{./logos.jpg}{}
}

\vspace*{5cm}

%{\centering
%\hspace*{0.5pt}\includegraphics[width=1\textwidth]{./ficha.pdf}{}}
}%
%\begin{minipage}{.8\textwidth}
%\centering\tiny\noindent{}Adverte-se aos curiosos que se imprimiu este livro \ifdef{\grafica}{na %gráfica \grafica}{em nossas oficinas}, 
%em \today \ifdef{\papelmiolo}{em papel \papelmiolo}, em tipologia \tipopadrao{}, com diversos %sofwares livres, 
%entre eles, Lua\LaTeX, git \& ruby. \ifdef{\RevisionInfo{}}{\par(v.\,\RevisionInfo)}{}\par \begin{%center}\normalsize\adforn{64}\end{center}
%\end{center}
